\subsection*{Chapter 6: Modern Architectures: Qwen3 from Scratch}

\paragraph{Exercise 6.1.}
If $\mu=0$ then $\sigma^2=\tfrac1d\sum_i(x_i-\mu)^2=\tfrac1d\sum_i x_i^2=\RMS(\vect x)^2$, so
$(\vect x-\mu)/\sqrt{\sigma^2+\epsilon}=\vect x/\sqrt{\RMS^2+\epsilon}$ and the outputs coincide.
In the explicit implementations the LayerNorm saves $x-\mu$, the denominator and the normalized
tensor (plus the input), while RMSNorm saves the fp32 input, the reciprocal RMS (one value per row)
and the normalized tensor; for a $(4,256,1024)$ fp32 input expect RMSNorm to save roughly one full
$4\cdot256\cdot1024\cdot4$\,B $=\SI{4.2}{\mega\byte}$ tensor less.

\paragraph{Exercise 6.2.}
$8d^2=3dh$ gives $h=\tfrac83d$. For $d=4096$ parity is $h\approx10923$; Llama~2-7B's $h=11008$ is
0.8\% larger. $11008=43\times256$: rounding to a multiple of 256 makes every GEMM tile full and lets
$h$ be split evenly across 2, 4 or 8 tensor-parallel ranks while keeping each shard a multiple of 32
or 64.

\paragraph{Exercise 6.3.}
Reorder the coordinates so that dimension $i$ and $i+\dhead/2$ are adjacent; in that basis RoPE at
position $m$ is the block-diagonal matrix $R_m=\mathrm{diag}(R(m\theta_0),\ldots,R(m\theta_{\dhead/2-1}))$.
Then $\langle R_mq,R_nk\rangle=q^\top R_m^\top R_nk=q^\top R_{n-m}k$ block by block, as in the 2-D
case; permuting both vectors identically does not change a dot product. Shifting $m$ and $n$ by the
same $c$ leaves $n-m$ unchanged. Left padding by $c$ tokens shifts every real token's position by $c$
within that row, so all scores between real tokens are unchanged (pad keys are masked out).

\paragraph{Exercise 6.4.}
With \code{p = halves\_to\_interleaved\_perm(D)}, the interleaved model's $q$ satisfies
$q_\text{inter}=q_\text{halves}[p]$ per head. Hence, for each head $h$, the rows
\code{W[h*D:(h+1)*D]} of the rotate-half checkpoint equal the interleaved rows reindexed by the
inverse permutation: \code{W\_halves[h*D + j] = W\_inter[h*D + inv(p)[j]]}; apply this to $\mat W_Q$
and $\mat W_K$ (and to the QK-norm gains if present). $\mat W_V$ is untouched: RoPE is not applied to
values, and the attention scores---the only place the permutation acts---are invariant.

\paragraph{Exercise 6.5.}
Qwen3-8B stores \SI{144}{\kibi\byte} per token in bf16, so \num{131072} tokens need
\SI{18}{\gibi\byte} ($\approx\SI{19.3}{\giga\byte}$); in fp8, \SI{9.7}{\giga\byte}. Next to
\SI{16.4}{\giga\byte} of bf16 weights a single context needs \SI{35.7}{\giga\byte}: it fits on an
80\,GB GPU, and with a 10\% reserve two such contexts fit ($\lfloor(72-16.4)/19.3\rfloor=2$). With
8-bit weights (\SI{8.2}{\giga\byte}) three fit with a bf16 cache and six with an fp8 cache.

\paragraph{Exercise 6.6.}
Per decode token and layer, attention costs $4TH\dhead$ FLOPs and reads $2TG\dhead b$ bytes of cache
($b$ bytes per element), so $I=2H/(Gb)$. In bf16 ($b=2$): $I=H/G=2,4,8$ for 0.6B, 8B, 32B; in fp8
($b=1$): $4,8,16$. Materializing \code{repeat\_interleave} writes and then reads an $H$-head copy, so
the bytes per FLOP grow by at least $H/G$ and the effective intensity falls back to about 1 (or below,
counting the extra write).

\paragraph{Exercise 6.7.}
If the cos/sin tables are computed in fp32 and only then cast to bf16, each entry has a relative
error of about $2^{-9}$, so the angle error stays near $0.004$ radians at every position. If, however,
the angles $m\theta_i$ themselves are formed in bf16 (e.g.\ because the model, including its buffers,
was cast before the tables were built), the fastest frequency has $\theta_0=1$ and the angle equals
$m$; bf16's spacing at magnitude $m\in[2^e,2^{e+1})$ is $2^{e-7}$, so the rounding error
($\le2^{e-8}$) exceeds $0.1$ rad once $e\ge5$, i.e.\ already at positions of a few dozen tokens.

\paragraph{Exercise 6.8.}
$\beta'=10^6\cdot4^{128/126}\approx\num{4.09e6}$. The frequency ratio for pair $i$ is
$\theta_i'/\theta_i=s^{-2i/(\dhead-2)}$, so the wavelength is multiplied by $s^{2i/(\dhead-2)}$:
unchanged for $i=0$ (smallest change) and multiplied by exactly $s=4$ for $i=\dhead/2-1=63$ (largest).

\paragraph{Exercise 6.9.}
Capacity $=\lceil c\,Nk/E\rceil=1.25\cdot4096\cdot2/64=160$ slots per expert. Under uniform random
routing an expert's load is approximately $\mathcal N(128,\,11.2^2)$. The expected overflow
$\E[(X-160)^+]=\sigma[\varphi(z)-z(1-\Phi(z))]$ with $z=2.85$ is $\approx0.007$ slots per expert,
i.e.\ $\approx0.006\%$ dropped. With $c=1.0$ ($z=0$) it is $\sigma\varphi(0)\approx4.5$ slots per
expert, about 3.5\% dropped; with $c=2.0$ ($z\approx11$) essentially none. Trained routers are far from
uniform, so real drop rates at $c=1$ are much higher, which is why dropless kernels became popular.

\paragraph{Exercise 6.10.}
Assume the \num{16384} tokens are spread evenly, 2048 per GPU. Each token sends $k=8$ copies of a
$d=2048$ bf16 vector, 7/8 of which leave the GPU: dispatch is $2048\cdot8\cdot2048\cdot2\cdot\tfrac78
\approx\SI{58.7}{\mega\byte}$ per GPU, and the combine step returns the same, \SI{117}{\mega\byte} per
layer, or about \SI{0.26}{\milli\second} at \SI{450}{\giga\byte\per\second}. Each GPU's experts process
about \num{16384} token-slots at $6dh=\SI{9.4}{M}$ FLOPs each, \SI{155}{GFLOP}, i.e.\
$\approx\SI{0.16}{\milli\second}$ at peak on an H100. Communication is comparable to (or larger than)
compute, so production systems overlap the all-to-all of one micro-batch with the expert GEMMs of
another.

\paragraph{Exercise 6.11.}
Naive loading of the 1.7B model in default fp32 holds the randomly initialized fp32 model
(\SI{6.9}{\giga\byte}) plus the bf16 state dict (\SI{3.4}{\giga\byte}); even in bf16 it holds two
copies (\SI{6.9}{\giga\byte}). The meta-device path peaks at one copy plus the largest single tensor,
the $\num{151936}\times2048$ embedding (\SI{0.62}{\giga\byte}): about \SI{4.1}{\giga\byte}.
